\chapter{总结与展望}
\label{chap:conclusion}

\section{工作总结}

本文围绕多轮对话检索中的\textbf{查询重写}问题，系统研究了"上下文依赖对话查询 $\rightarrow$ 语义自包含独立查询"这一关键预处理模块的建模、训练与对齐问题，提出并实现了由\textbf{历史剪枝—渐进式重写—检索反馈对齐}三阶段构成的查询重写框架 \textsc{PruneRewriteAlign}（PRA）。本文的主要工作和创新点如下：

（1）\textbf{话题漂移感知的对话历史剪枝（HP-Pruner）}。针对现有方法"把整段历史一视同仁地拼接"带来的噪声注入、计算浪费与话题切换误导问题，本文首次在轮次级引入三分类相关性判别器，以 ModernBERT-base 作为轻量骨干，基于 TopiOCQA 话题切换标注、QReCC 人工重写对比构造了 94\,422 条弱监督训练样本。在 4\,721 条留出验证集上 HP-Pruner 取得 80.1\% 准确率与 0.797 的宏 F1，为后续重写器提供了干净、紧凑的历史上下文。

（2）\textbf{两阶段渐进式查询重写（PRW）}。针对指代消解与省略补全在语言学性质、错误模式上的显著异质，本文提出把端到端重写任务解耦为"Stage-1 仅指代消解—Stage-2 在此基础上补全省略"的串行结构。基于 Qwen2.5-3B-Instruct 与 LoRA（$r{=}16$），PRW 在 TopiOCQA、Doc2Dial 这类话题切换密集、上下文依赖强的基准上将 Recall@5 相对 No-Rewrite 分别提升 39 与 66 个百分点以上（BGE-large 检索器下），在 3 种异构检索器上表现一致。

（3）\textbf{检索反馈驱动的偏好对齐（RFR）}。针对"像人写的参考重写 $\ne$ 被检索器召回的好重写"这一目标失配问题，本文提出将 BGE-large 对"候选重写 vs gold passage"的余弦相似度作为偏好信号、以 DPO 进行轻量化对齐的 RFR 方案。从 10\,000 个 TopiOCQA 训练轮次采样 80\,000 候选后构造 4\,106 条高质量偏好对，DPO 训练在 129 步内把 reward accuracies 从 45\% 稳定提升到 85\%。在 4 基准 $\times$ 3 检索器共 12 组设定中，RFR 相对其 PRW 起点取得平均 $+1.10$ R@5 / $+1.31$ R@20 的稳定增益，且该增益可迁移到训练从未见过的 GTE-ModernBERT / Dragon-Multiturn 检索器。

（4）\textbf{全面的实证分析与案例研究}。本文在 QReCC、QuAC、Doc2Dial、TopiOCQA 4 个公开基准上真实跑通了 HP-Pruner $\to$ PRW $\to$ RFR 的完整链路；并在实验室合作构造的 MTR-Bench 外部评测集上验证了其标注信度（Fleiss $\kappa\in[0.58, 0.78]$、Krippendorff $\alpha\in[0.71, 0.90]$，达到 "substantial/almost-perfect" 一致性），确认其可作为本文方法在极端场景下的可信外部参考。

综合来看，本文从"噪声过滤—结构化重写—下游反馈"三个环节出发，构建了一套低成本、高质量、可落地的多轮对话查询重写范式，对工业级对话 RAG 系统具有直接参考价值。

\section{局限性}

尽管本文方法取得了一致性的性能提升，但仍存在以下若干局限性：

\textbf{语言与领域局限}：本文实验主要在英文通用领域开展，尚未系统验证方法在中文对话、多语言混合、垂直专业领域（医疗、法律、金融）上的迁移能力。

\textbf{检索器假设}：RFR 的偏好信号由单个稠密检索器提供，尚未考虑"检索—重排—生成"全链路的联合反馈。

\textbf{合成数据的噪声}：PRW 第一阶段的中间参考由自动脚本构造，虽经 GPT-4o 校验但仍存在少量噪声，可能对小规模模型的稳定性产生限制。

\textbf{对极端长对话的鲁棒性}：现有实验的对话轮数普遍小于 20 轮；在更长的用户历史场景下（如长期记忆对话助手），HP-Pruner 的分类器可能需要进一步扩展至跨段剪枝。

\textbf{QuAC 场景的回退}：在 QuAC 这类"答案即前文窗口"的伪多轮基准上，由于历史答案段直接匹配是 No-Rewrite 基线的天然优势，PRW/RFR 的重写反而会打偏这一方向；后续可在场景识别层面对此类基准做自适应 bypass。

\section{未来展望}

基于本文工作可扩展的研究方向包括：

\textbf{多信号融合的反馈对齐}：将生成端的答案忠实度、用户满意度等信号与检索器 Recall 联合作为 DPO 的偏好信号，形成更加全链路对齐的训练目标。

\textbf{端到端对话检索与重写的联合优化}：探索重写器与检索器的联合蒸馏或联合训练，在保留重写可解释性的同时进一步提升端到端性能。

\textbf{面向长时记忆与跨会话的剪枝}：将 HP-Pruner 扩展至"跨会话、跨天"的长时对话场景，结合记忆压缩与摘要机制。

\textbf{低资源语言的迁移}：通过跨语言对齐、弱监督数据生成等手段把 PRA 框架推广至中文等低资源对话检索场景。

\textbf{与生成式检索器的结合}：探索本文方法在生成式检索（如 DSI、GENRE）架构下的适配与优化，推动下一代对话检索系统的演进。

\section{结束语}

对话式人工智能正处于从"单轮搜索 / 问答"向"多轮理解 / 协作"深度演进的关键阶段。查询重写作为连接用户自然语言与底层检索系统的关键桥梁，其重要性将随着应用场景的复杂化而持续上升。本文所做的努力只是这条漫长道路上的一小步，期望本文提出的框架与分析能够为学术界与工业界研发更加轻量、更加稳健、更加对齐的对话式 RAG 系统提供一点参考与启发。
